\def\topic{Induction \& Recursion}
\ifdefined\MyHeaderIncluded
\else
\def\qwertyflag{true}

\ifdefined\MyHeaderIncluded

\else
  \newcommand{\foobar}{FooBar}
  \def\MyHeaderIncluded{true}
% if is till the end of the file

\documentclass[10pt,a4paper]{article}
\usepackage[utf8]{inputenc}
%\usepackage[german]{babel}
%\usepackage[cyrillic]{babel}
\usepackage[T1]{fontenc}
%\usepackage[T2A]{fontenc}
%\usepackage[english, russian]{babel}
\usepackage{amsmath}
\usepackage{amsfonts}
\usepackage{amssymb}
\usepackage{amsthm}
\usepackage{mathtools} %coloneqq

\usepackage{bussproofs} %

%\renewcommand{\thesubsection}{\arabic{subsection}}
%\newtheorem{thm}{Theorem}[section]
%\renewcommand{\thethm}{\arabic{subsection}.\arabic{thm}}

\usepackage{thmtools}
\usepackage{hyperref}
\hypersetup{
    colorlinks=true,
    linkcolor=blue,
    filecolor=magenta,      
    urlcolor=cyan,
    pdftitle={Overleaf Example},
    %pdfpagemode=FullScreen,
    }
%\usepackage{stix-type1}%\turnediota
\theoremstyle{definition}

\newtheorem{constructor}{Constructor}[section]
\newtheorem{iaxiom}{Axiom}
\newtheorem{axiom}{Axiom}
\newtheorem*{axiom*}{Axiom}
\newtheorem{definition}{Definition}[section]
\newtheorem*{definition*}{Definition}
\newtheorem{remark}{Remark}[section]
\newtheorem*{remark*}{Remark}
\newtheorem{recall}{Recall}[section]
\newtheorem*{recall*}{Recall}

\newtheorem{lemma}{Lemma}[section]
\newtheorem{theorem}{Theorem}[section]
\newtheorem*{theorem*}{Theorem}
\newtheorem{note}{Note}[section]
\usepackage{etoolbox}

%\newcommand*{\@Dashv}[2]{%
%  \reflectbox{$\m@th#1#2$}%
%}
\newcommand*{\separation}[3]{% 1)binder, 2)set, 3)condition
  \{#1\in #2: #3\}
}

% https://www.freecodecamp.org/news/simple-regex-tricks-for-beginners-3acb3fa257cb/
%  \\mathcal\{T\} |---> \topology
%  \\mathcal\{([^\}]*)\}   |--->  $1
\newcommand{\AND}{\ \ \&\ \ }
\newcommand{\any}{\mathrm{\ any}}
\newcommand{\borel}{\mathcal{B}}
\newcommand{\can}{\mathrm{can}}
\newcommand{\canonize}{\mathrm{canonize}}
\newcommand{\closure}{\mbox{cl}}
\newcommand{\complex}{\mathbb{C}}
\newcommand{\converges}{\longrightarrow}
\newcommand{\countable}{\mbox{\ countable}}
\newcommand{\ldefined}{\mbox{\ defined}}
\newcommand{\dom}{\mbox{dom}}
\newcommand{\Dom}{\mbox{Dom}}
\newcommand{\Done}{\noindent \textbf{Done.}\\}
\newcommand{\lexists}{\mbox{\ exists}}
\newcommand{\lforall}{\mbox{\ forall}}
\newcommand{\lunique}{\mbox{\ unique}}
\newcommand{\family}{\mathrm{\ family}}
\newcommand{\Fm}{\mathrm{Fm}}
\newcommand{\FuncSymb}{\mathrm{FuncSymb}}
\newcommand{\function}{\mathrm{\ function}}
\newcommand{\functional}{\mbox{\ functional}}
\newcommand{\IFF}{\ \Longleftrightarrow\ }
\renewcommand{\iff}{\ \longleftrightarrow\ }
\newcommand{\IMPLIES}{\ \Longrightarrow\ }
\renewcommand{\implies}{\ \longrightarrow\ }
\newcommand{\Inf}{\mathrm{Inf}}
\newcommand{\injective}{\mbox{\ injective}}
\newcommand{\injection}{\mbox{\ injection}}
\newcommand{\interior}{\mbox{int}}
\newcommand{\intersection}{\bigcap}
\newcommand{\isempty}{\mathrm{\ empty}}
\newcommand{\LB}{\mathrm{LB}}
\newcommand{\nghbset}{\mathcal{N}}
\newcommand{\nonempty}{\mathrm{\ nonempty}}
\newcommand{\Ord}{\mathrm{Ord}}
\newcommand{\power}{\mathcal{P}}
\newcommand{\pr}{\mathrm{pr}}
\newcommand{\fst}{\mathrm{fst}}
\newcommand{\snd}{\mathrm{snd}}
\newcommand{\PredSymb}{\mathrm{PredSymb}}
\newcommand{\Proof}{\noindent \textbf{Proof.}\\}
\newcommand{\ran}{\mathrm{ran}}
\newcommand{\Ran}{\mathrm{Ran}}
\newcommand{\real}{\mathbb{R}}
\newcommand{\relation}{\mathrm{\ relation}}
\newcommand{\set}{\mathrm{\ set}}
\newcommand{\singleton}{\mathrm{\ singleton}}
\newcommand{\semantics}{\mathrm{semantics}}
\newcommand{\semialgebra}{\mathcal{S}}
\newcommand{\Sup}{\mathrm{Sup}}
\newcommand{\surjective}{\mbox{\ surjective\ on\ }}
\newcommand{\Term}{\mathrm{Term}}
\newcommand{\Type}{\mathrm{Type}}
\newcommand{\topology}{\mathcal{T}}
\newcommand{\total}{\mbox{\ total\ on\ }}
\newcommand{\transitive}{\mathrm{\ transitive}}
\newcommand{\union}{\bigcup}
\newcommand{\lTrue}{\mathrm{True}}
\newcommand{\lFalse}{\mathrm{False}}
\newcommand{\lUncertain}{\mathrm{Uncertain}}
\newcommand{\lif}{\mathrm{if}}
\newcommand{\lthen}{\mathrm{then}}
\newcommand{\lelse}{\mathrm{else}}
\newcommand{\lexception}{\mathrm{exception}}
\newcommand{\lundef}{\mathrm{undef}}
\newcommand{\univocal}{\mathrm{\ univocal}}
\newcommand{\univocalat}{\mathrm{\ univocal\ at\ }}
\newcommand{\UB}{\mathrm{UB}}
\newcommand{\Un}{\mathcal{U}}
\newcommand{\Var}{\mathrm{Var}}

\newcommand{\notvarphi}{\phi}
\renewcommand{\phi}{\varphi}

\newcommand{\comment}[1]{\quad\mbox{(#1)}}
% local commands
%\preto\section{\undef\blipo} %etoolbox
\preto\subsection{\undef\blipo} %etoolbox
% https://tex.stackexchange.com/questions/105684/limit-macro-newcommand-scope-to-section
\usepackage{graphicx}
\usepackage{authblk}

\setlength\parindent{0pt} % whole doc of \noindent
\usepackage{soul}
%\usepackage{mathgl}

\fi % guard
\ifdefined\topic
\else
\def\topic{No Titel}
\fi
\title{\topic}
\author[1]{Georgy Dunaev}
\begin{document}
\maketitle
\fi
\section{Height of the set}
%Theorems about transifinite- and epsilon- induction and recursion.
\begin{recall}\ \\
Transfinite recursion\\
$G:\V\to\V$\\
$\exists! F:\Ord\to\V$, s.t.\\
$F(a) = G(F\restriction a)$.
\end{recall}
\begin{definition}[union]\ \\
\[
\bigcup^y_{\psi} t 
\ \eqdef\ 
\cc{x}{\exists y, \psi \land x\in t}
\]
Side conditions:\\
1) formula $\psi$ must not depend on variable $x$,\\
2) term $t$ must not depend on variable $x$,\\
3) variable $y$ is different from variable $x$.
%Side conditions:\\
%1) variable $x$ is not free in formula $\psi$,\\
%2) variable $x$ is not free in term $t$,\\
%3) variable $x$ is not free in variable $y$.\\
%Dependencies:\\
%1) formula $\psi$ may depend on $y$.\\
%2) term $t$ may depend on $y$.
\end{definition}

\begin{definition}[intersection of a definable conglomerate of classes]\ \\
\[
\bigcap^y_{\psi} t 
\ \eqdef\ 
\cc{x}{\forall y, \psi \implies x\in t}
\]
Side conditions:\\
1) formula $\psi$ must not depend on variable $x$,\\
2) term $t$ must not depend on variable $x$,\\
3) variable $y$ is different from variable $x$.
\end{definition}
Variable is a syntatic atom.\\
Object is a semantic atom.\\

Functional is a syntactic function.\\
Family is a semantic function.\\

Evaluation is a syntactic image.
\begin{notation}[intersection of evaluations of the elements of a class]\ \\
Let $y$ be a variable, let $s$ and $t$ be terms. Then for the case if formula has form $(y\in s)$, and if $y$ is also its bound variale, then it is allowed to skip writing a variable on top.
\[
\bigcap_{y\in s} t 
\ \eqdef\ 
\bigcap^y_{y\in s} t 
\]
\end{notation}
\begin{notation}[intersection of the elements of a class]\ \\
\[
\bigcap G = \cc{x}{\forall y\in G. x\in y}
\]
\[
\bigcap G 
\ \eqdef\ 
\bigcap_{y\in G} y 
\]
\end{notation}
\begin{definition}[W]\ \\
The class of well-founded sets $\W$ is being defined as an intersection of all those classes, which contain each of their subsets in themselves.
\[
\W\eqdef\cc{x}{\forall Q, (\power Q\subclass Q)\implies (x\in Q)}
\]
\end{definition}
\begin{notation}
For such intersections of proper classes satisfying some property the following notation may be used:
%\[
%\W\eqdef \bigwedge_Q (\power Q\subclass Q).
%\]
\[
\W\eqdef \bigcap^Q\limits_{(\power Q\subclass Q)} Q.
\]

\[
\bigcap_{\psi(Q)} t(Q) 
\eqdef
\cc{x}{\forall Q, \psi(Q) \implies x\in t(Q)}
\]



\[
\bigcap_{y\in Y} t(y)
=\bigcap\cc{t(y)}{y\in Y}
=\cc{m}{\forall y\in Y, m\in t(y)}
=\cc{m}{\forall y, y\in Y \implies m\in t(y)}
\]
\[
=\cc{x}{\forall Q, Q\in Y \implies x\in t(Q)}
\]

\[
\bigcap_{y\in Y} t(y)
= \bigcap \cc{t(y)}{{y\in Y}}
= \bigcap (y\mapsto t(y))[Y]
\]
\end{notation}
\begin{recall}\ \\
Epsilon recursion:\\
$G:V\to V$\\
$\exists! F:\W\to\V$, s.t.\\
$F(a) = G(F\restriction a)$.
\end{recall}
\begin{lemma}[W]\ \\
\[
\forall Q, (\power Q\subclass Q) \implies (\W\subclass Q)
\]
\Proof
Fix $Q$. Assume $(\power Q\subclass Q)$. To Show: $\W\subclass Q$.\\
%Proposition: $\W\subclass Q$.\\
Fix $m$.\\
Assume $m\in\W$.\\
To show: $m\in Q$.
\[
(m\in\W)
\IMPLIES
\forall D, (\power D\subclass D)\implies (x\in D)
\]
Set $D:=Q$.\\
The assumption $(\power Q\subclass Q)$ together with the last statement lead to $(m\in Q)$ via specification and modus ponens.
%So the statement $(\W\subclass Q)$ is proven.
\DoneNoPar
\end{lemma}
\begin{note}
Due to this lemma, the axiom of regularity may be formulated as $(\V\subclass\W)$.
\end{note}
\begin{theorem}[W]\ \\
The class $\W$ contains all of its subsets in itself.
\[
\power \W \subclass \W
\]
\Proof
Fix $(x\in \power \W)$. To show: $x\in\W$.
\[
(x\in \power \W)
\IMPLIES
(x\set)
\land
(x\subclass\W)
\]
To show: $(x\set)\land(\forall Q, (\power Q\subclass Q)\implies (x\in Q))$.\\
Fix $Q$.\\
Assume $(\power Q\subclass Q)$.\\
To show $(x\in Q)$.\\
The propostion $(\W\subclass Q)$ is proven by the specification of the previous lemma and then by modus ponens of the result together with the assumption.
\[
(x\subclass\W)
\LAND
(\W\subclass Q)
\IMPLIES
(x\subclass Q)
\]
\[
(x\set)
\LAND
(x\subclass Q)
\IMPLIES
(x\in\power Q)
\]

\[
(x\in\power Q)
\LAND
(\power Q\subclass Q)
\IMPLIES
(x\in Q)
\]

\Done
\end{theorem}


\begin{definition}
Height function definition:\\
Let $h(x)$ denote the height of a the set, which is the maximum length of a $\ni$-chain, started from $x$, minus one.\\
\end{definition}

\begin{example*}\ \\
$h(\emptyclass)$=0,\\
$h(\{\emptyclass\})$=1,\\
$h(\{\emptyclass,\{\emptyclass\}\})$=2,\\
and in general if $n$ is a finite ordinal number, then $h(n)=n$.
\end{example*}
\begin{definition}[height]\ \\
\label{height}
The height function is defined via $\in$-recursion:
\[
h(x) \ceq \sup \cc{ h(y) }{ y \in x}.
\]
\end{definition}
Property:
\[
h(x) = \sup h[x]
\]

$G(R) \ceq \sup(\ran(R)) = \union(\ran(R))$\\
via epsilon recursion.\\
$h : \V\to\V$\\
$h(a) = \sup(\ran(h\restriction a) )$\\


% height = depth
\begin{lemma}[set's height is ordinal]\ \\
\label{lem:depthinon}
The height of any set is an ordinal.
\[
(h:V\to\On)
\]
%\omega
\end{lemma}
\ProofNoPar\ (first variant)\\
Using the epsilon induction.\\
Base of the epsilon induction, which may be ommited:
\[
h(\emptyclass) = \sup_\N h[\emptyclass] = \sup_\N \emptyclass = 0
\]


Step of the epsilon induction:\\
Assume $\forall y\in x: h(y)\in \On$.\\
To show: $h(x)\in\On$.\\
It follows from the assupmtion, that $h[x]\subclass\On$.\\
\[
\forall B, \left(B\subclass\On\middle) \implies \middle(\Ord\middle(\union B\middle)\right)
\]
The $x$ is a set; $h$ with a domain restricted to x is a function.\\
Therefore $h[x]$ is also a set.\\
Therefore $\union h[x]$ is also a set.\\

\[
\left(\union h[x]\set\middle)
\LAND
\Ord\middle(\union h[x]\middle)
\IMPLIES 
\middle(\union h[x]\in\On\right)
\]
\[
h(x) = \sup h[x] = \union h[x]
\]
Therefore $(h(x)\in\On)$.
\Done

\ProofNoPar\ (second variant)\\
Using the epsilon induction.\\
\[
\power Q\subclass Q \implies \V\subclass Q
\]
\[
Q\eqdef \cc{x\in\V}{h(x)\in\On}
\]
If $\V\subclass Q$ then $\forall x\in\V, h(x)\in\On$, and the proof is done.\\
To Show: $\power Q\subclass Q$, or, equivalently, $\forall x, (x\set)\land (x\subclass Q) \implies (x\in Q)$.\\
\[
(x\subclass Q)
\IMPLIES
(h[x]\subclass \On).
\]
\[
\IMPLIES
\Ord(\union h[x])
\]
\[
\IMPLIES
\Ord(\sup h[x])
\]
\[
\IMPLIES
\Ord(h(x))
\]
...

\[
(s\set)
\LAND
(h\mathrm{\ set\-valued\ function})
\IMPLIES
(h(x)\set)
\]
\[
(h(x)\set)
\LAND
\Ord(h(x))
\IMPLIES
(h(x)\in\On)
\]
\[
(x\in\V)
\LAND
(h(x)\in\On)
\IMPLIES
(x\in Q)
\]
\Done


\ifdefined\qwertyflag
\undef\qwertyflag
\end{document}
\fi